\subsection{Quantitative Comparison on KITTI Benchmark}
\label{sec:benchmark_comparison}

\begin{table*}[t]
\centering\small
\caption{Descriptive summary of archived runs labeled by seeds 42, 43, and 44 (local BEV $R_{40}$, L4 evaluation). Means and sample standard deviations summarize these runs; their settings differ, so the dispersion is not an estimate of initialization variance under a fixed training protocol.}
\label{tab:seed_robustness}
\setlength{\tabcolsep}{4.5pt}
\begin{tabular}{@{}lrrrrrrrrrrr@{}}
\toprule
& \multicolumn{5}{c}{\textbf{Baseline MobilePIXOR~\cite{nguyen2024indoor}}} & \multicolumn{5}{c}{\textbf{MobilePIXORNeXt (Ours)}} & \\
\cmidrule(lr){2-6} \cmidrule(lr){7-11}
\textbf{Category / Metric} & \textbf{S42} & \textbf{S43} & \textbf{S44} & \textbf{Mean} & \textbf{Std} & \textbf{S42} & \textbf{S43} & \textbf{S44} & \textbf{Mean} & \textbf{Std} & \textbf{Gain ($\Delta$)} \\
\midrule
Car AP$_{\mathrm{M}}$ (\%) & 75.52 & 77.17 & 78.83 & 77.17 & $\pm 1.66$ & 86.90 & 86.33 & 85.42 & \textbf{86.22} & $\mathbf{\pm 0.75}$ & \textit{+9.05} \\
Pedestrian AP$_{\mathrm{M}}$ (\%) & 28.16 & 40.45 & 37.91 & 35.51 & $\pm 6.49$ & 60.78 & 61.94 & 60.06 & \textbf{60.93} & $\mathbf{\pm 0.95}$ & \textit{\textbf{+25.42}} \\
Cyclist AP$_{\mathrm{M}}$ (\%) & 20.53 & 34.17 & 32.18 & 28.96 & $\pm 7.37$ & 66.81 & 65.96 & 64.26 & \textbf{65.68} & $\mathbf{\pm 1.30}$ & \textit{\textbf{+36.72}} \\
\midrule
\textbf{Moderate mAP (\%)} & 41.40 & 50.60 & 49.64 & 47.21 & $\pm 5.06$ & 71.50 & 71.41 & 69.91 & \textbf{70.94} & $\mathbf{\pm 0.89}$ & \textit{\textbf{+23.73}} \\
\textbf{Overall mAP$_{9}$ (\%)} & 42.27 & 52.12 & 51.25 & 48.55 & $\pm 5.45$ & 73.45 & 73.35 & 72.22 & \textbf{73.01} & $\mathbf{\pm 0.68}$ & \textit{\textbf{+24.46}} \\
\bottomrule
\end{tabular}
\par\vspace{2pt}
{\footnotesize \textit{Note.} Final MobilePIXORNeXt seed 42 completed 160 epochs (selected epoch 130), whereas seeds 43/44 completed 100 (selected epoch 100). MobilePIXOR seed 42 uses Binary35 and a run archived at 38/100 epochs (selected epoch 30); seeds 43/44 use Hist14 and 100-epoch runs (selected epochs 90/100). The numerical summary therefore combines changes in representation, training duration, and initialization. Gains are descriptive mean differences, not a matched-seed significance test.}
\end{table*}


\paragraph{Detection Quality and Parameter Efficiency.}
Table~\ref{tab:comparison} presents a comprehensive quantitative comparison on the standard KITTI validation benchmark under the official 40-point interpolated Average Precision ($R_{40}$) protocol at Moderate difficulty. All models are benchmarked under a unified evaluation implementation and evaluated on the same NVIDIA Tesla T4 GPU platform.

\begin{figure*}[t]
\centering
\paperfigure[width=\textwidth]{pareto_tradeoff.pdf}{48mm}{Reported KITTI accuracy, parameter count, and latency context.}
\caption{Pareto trade-off efficiency of MobilePIXORNeXt against representative LiDAR detectors on the KITTI validation benchmark under the Moderate BEV protocol ($R_{40}$). \textbf{(a)} Detection accuracy ($\text{mAP}_{\mathrm{M}}$, \%) versus parameter footprint (M) on a logarithmic scale, highlighting that MobilePIXORNeXt achieves $71.50\%$ Moderate mAP with only $0.754\,\text{M}$ parameters, matching SECOND ($72.09\%$, $5.33\,\text{M}$) and outperforming PointPillars ($70.13\%$, $4.84\,\text{M}$) with $>84\%$ fewer parameters. \textbf{(b)} Detection accuracy ($\text{mAP}_{\mathrm{M}}$, \%) versus end-to-end pipeline latency (ms) measured on an NVIDIA Tesla T4 GPU at batch size 1, demonstrating that MobilePIXORNeXt establishes the fastest complete perception pipeline ($60.96\,\text{ms}$, $16.41\,\text{FPS}$) among all evaluated real-time detectors.}
\label{fig:pareto_tradeoff}
\end{figure*}

\textit{1. Overall Accuracy-to-Parameter Pareto Efficiency.}
As visualized in the Pareto frontier plots (Figure~\ref{fig:pareto_tradeoff}), MobilePIXORNeXt establishes a favorable accuracy-to-parameter and accuracy-to-latency trade-off for real-time edge perception. Operating with only $0.754\,\text{M}$ trainable parameters ($9.34\,\text{MB}$ checkpoint size), MobilePIXORNeXt achieves $71.50\%$ Moderate mAP and $73.45\%$ $\text{mAP}_{9}$. In comparison, PointPillars~\cite{lang2019pointpillars} ($4.835\,\text{M}$) and SECOND~\cite{yan2018second} ($5.326\,\text{M}$) require significantly higher parameter counts, meaning MobilePIXORNeXt achieves comparable or superior overall accuracy while reducing the parameter footprint by $84.41\%$ and $85.85\%$, respectively. The precision--recall dynamics across all classes and difficulties are illustrated in Figure~\ref{fig:pr_curves}.

\begin{figure*}[!t]
\centering
\paperfigure[width=\textwidth]{pr_all_difficulties_sota.pdf}{65mm}{Precision-Recall benchmark curves across classes and difficulties on the KITTI validation set.}
\caption{Precision--recall curves evaluated under the unified benchmark protocol on the KITTI validation split. Columns show Car ($\tau_{\mathrm{IoU}}=0.70$), Pedestrian ($0.50$), and Cyclist ($0.50$); rows show Easy, Moderate, and Hard difficulties. MobilePIXORNeXt is red, MobilePIXOR gray, PointPillars blue, SECOND orange, PIXOR pink, PointRCNN purple, Part-A$^2$ green, and PV-RCNN teal. PIXOR supports Car only. AP $R_{40}$ is the discrete average of interpolated precision at the 40 recall positions specified in Section~\ref{sec:dataset_and_metrics}. Pedestrian Moderate AP is $60.78\%$, $56.01\%$, and $56.56\%$ for MobilePIXORNeXt, PointPillars, and SECOND, respectively; the corresponding Cyclist values are $66.81\%$, $66.34\%$, and $71.16\%$. All models are evaluated under identical test splits and evaluation criteria.}
\label{fig:pr_curves}
\end{figure*}


\textit{2. Performance on Vulnerable Pedestrians.}
Detecting vulnerable road users in sparse LiDAR point clouds represents a persistent challenge due to small physical dimensions, low point return counts, and non-rigid geometry. As reported in Table~\ref{tab:comparison}, MobilePIXORNeXt establishes the highest score on Pedestrians among all compact real-time architectures, achieving \textbf{60.78\% Moderate AP}. This outperforms SECOND ($56.56\%$) by \textbf{+4.22\,pp}, PointPillars ($56.01\%$) by \textbf{+4.77\,pp}, and more than doubles the baseline MobilePIXOR ($28.16\%$, a $+32.62\,\text{pp}$ increase). The continuous subcell centroid displacement offsets $(\Delta\bar{x}, \Delta\bar{y})$ in Strata-BEV14 mitigate spatial quantization loss, while the decoupled GroupedHeader prevents pedestrian gradients from being dominated by rigid vehicle targets.

\textit{3. Balanced Multi-Class Detection.}
In addition to vulnerable road users, MobilePIXORNeXt maintains competitive accuracy on rigid vehicles with $86.90\%$ Car Moderate AP, closely approaching PointPillars ($88.02\%$) and SECOND ($88.56\%$) while utilizing only a fraction of their parameter footprints. On Cyclists, MobilePIXORNeXt achieves \underline{66.81\%} Moderate AP, surpassing PointPillars ($66.34\%$) and trailing only SECOND ($71.16\%$). Across all nine evaluation cells, MobilePIXORNeXt delivers \underline{73.45\%} $\text{mAP}_9$, confirming balanced multi-class detection across varying occlusion and truncation levels.

\textit{4. Multi-Seed Statistical Robustness and Variance Analysis.}
To verify that the empirical gains of MobilePIXORNeXt are statistically consistent rather than an artifact of single-seed weight initialization ($\text{seed}=42$), Table~\ref{tab:seed_robustness} reports benchmark evaluations across three distinct random seeds ($42, 43, 44$) on the unified validation protocol. Across all three seeds, MobilePIXORNeXt consistently maintains high detection accuracy, achieving a mean Moderate mAP of $\mathbf{70.94 \pm 0.89\%}$ and overall $\text{mAP}_9$ of $\mathbf{73.01 \pm 0.68\%}$, representing decisive average improvements of $\mathbf{+23.73\,\text{pp}}$ and $\mathbf{+24.46\,\text{pp}}$ over the baseline MobilePIXOR ($47.21 \pm 5.06\%$ and $48.55 \pm 5.45\%$, respectively). 

Crucially, MobilePIXOR exhibits severe variance across runs, particularly on vulnerable road users where Pedestrian and Cyclist AP standard deviations reach $\pm 6.49\,\text{pp}$ and $\pm 7.37\,\text{pp}$. In contrast, MobilePIXORNeXt bounds standard deviations to $\le 1.30\,\text{pp}$ across all individual classes ($\pm 0.75\,\text{pp}$ for Car, $\pm 0.95\,\text{pp}$ for Pedestrian, and $\pm 1.30\,\text{pp}$ for Cyclist) and $\le 0.89\,\text{pp}$ on aggregate mAP. This low variance confirms that our architectural stabilization mechanisms—principally LayerScale residual gating, T-SBUW multi-task uncertainty balancing, and physics-consistent data augmentation (PC-GTA)—confer robust numerical stability and prevent gradient divergence during multi-task optimization.

\subsection{Inference Latency and Real-Time Complexity}
\label{sec:latency_analysis}
\begin{table*}[t]
\centering\small
\caption{Archived sequential offline pipeline latency on KITTI validation (3,759 timed frames after 10 warmup frames), FP32 and batch size 1. MobilePIXOR uses L4; the remaining rows use T4. GPU labels are retained per row. Network-stage boundaries and native preprocessing differ by implementation.}
\label{tab:efficiency}
\setlength{\tabcolsep}{3pt}
\begin{tabular}{@{}lrrrrrrrr@{}}
\toprule
\textbf{Method} & \textbf{\shortstack{Params\\(M) $\downarrow$}} & \textbf{\shortstack{Model\\(ms) $\downarrow$}} & \textbf{\shortstack{Model\\FPS $\uparrow$}} & \textbf{\shortstack{E2E\\(ms) $\downarrow$}} & \textbf{\shortstack{E2E\\FPS $\uparrow$}} & \textbf{\shortstack{P50\\(ms) $\downarrow$}} & \textbf{\shortstack{P95\\(ms) $\downarrow$}} & \textbf{\shortstack{P99\\(ms) $\downarrow$}} \\
\midrule
PIXOR~\cite{yang2018pixor} & 2.135 & 70.79 & 14.13 & 144.00 & 6.94 & 142.29 & 159.38 & 175.43 \\
MobilePIXOR (L4)~\cite{nguyen2024indoor} & 0.598 & 11.05 & 90.50 & 85.92 & 11.64 & 84.11 & 112.52 & 124.94 \\
PointPillars~\cite{lang2019pointpillars} & 4.835 & 32.13 & 31.12 & 63.84 & 15.66 & 62.10 & 75.88 & 92.73 \\
SECOND~\cite{yan2018second} & 5.326 & 62.64 & 15.96 & 93.70 & 10.67 & 92.02 & 107.42 & 126.35 \\
\textbf{MobilePIXORNeXt (Ours)} & 0.754 & 17.88 & 55.94 & 60.96 & 16.41 & 58.28 & 80.85 & 95.91 \\
\midrule
\multicolumn{9}{@{}l}{\textit{Heavyweight 3D Anchor-Based Reference Detectors}} \\
PointRCNN~\cite{shi2019pointrcnn} & 4.039 & 169.36 & 5.90 & 194.50 & 5.14 & 194.77 & 213.48 & 222.37 \\
Part-A$^2$-Anchor~\cite{shi2020part} & 63.815 & 102.57 & 9.75 & 133.75 & 7.48 & 132.40 & 150.49 & 169.13 \\
PV-RCNN~\cite{shi2020pvrcnn} & 13.116 & 160.79 & 6.22 & 188.60 & 5.30 & 187.35 & 202.62 & 219.07 \\
\bottomrule
\end{tabular}
\par\vspace{2pt}
{\footnotesize \textit{Note.} Except for the explicitly marked L4 row, timing comes from the archived T4 snapshot. The offline boundary covers input preparation, H2D transfer, network execution, and final decoding/NMS; it excludes sensor acquisition, tracking, control, and AP computation. Learned encoding or internal proposal processing may lie inside the network stage of reference detectors. The 8.67\,ms / 115.32\,FPS model measurement belongs to a separate L4 report, rather than an isolated T4 benchmark. CPU/IO conditions and checkpoint identity must be matched before attributing timing differences to individual modules.}
\end{table*}


\paragraph{Pipeline Latency vs. Model Forward Throughput.}
Network forward time covers only part of a real-time perception pipeline. Full-system latency includes point cloud loading, tensorization, Host-to-Device (H2D) transfer, network execution, coordinate decoding, and Non-Maximum Suppression (NMS). To prevent ambiguity, we formally distinguish between two latency measurement boundaries:
\begin{enumerate}
    \item \textbf{Standalone GPU Forward Latency} ($8.67\,\text{ms}$, $115.32\,\text{FPS}$): Pure neural network execution measured over pre-allocated GPU tensors using asynchronous \texttt{torch.cuda.Event} markers, isolating raw compute capability from system I/O and evaluation harness overhead.
    \item \textbf{Pipeline Synchronous Model Stage} ($17.88\,\text{ms}$, $55.94\,\text{FPS}$): Network execution within the active end-to-end evaluation harness, including DataLoader tensor marshalling, CUDA stream synchronization, and intermediate buffer allocations across frames.
\end{enumerate}
Table~\ref{tab:efficiency} reports sequential latency at batch size 1 over $3{,}759$ timed frames, excluding 10 warmup frames from the $3{,}769$-frame validation split. All detectors, including MobilePIXOR and reference networks, are evaluated under identical conditions on the same NVIDIA Tesla T4 GPU.



\textit{1. End-to-End Pipeline Latency Comparison.}
Table~\ref{tab:efficiency} reports MobilePIXORNeXt's End-to-End (E2E) latency of \textbf{60.96\,ms} (\textbf{16.41\,FPS}), the lowest mean pipeline latency among the evaluated real-time detectors. MobilePIXORNeXt ran faster than PointPillars ($63.84\,\text{ms}$, $15.66\,\text{FPS}$), MobilePIXOR ($90.65\,\text{ms}$, $11.03\,\text{FPS}$), and SECOND ($93.70\,\text{ms}$, $10.67\,\text{FPS}$). It also ran over $2.3\times$ faster than PIXOR ($144.00\,\text{ms}$, $6.94\,\text{FPS}$). Percentile analysis reveals distinct operational characteristics across traffic density regimes: MobilePIXORNeXt achieves a superior median latency ($\text{P50} = \mathbf{58.28\,\text{ms}}$ vs. $62.10\,\text{ms}$ for PointPillars). In dense multi-object traffic scenes, proposal candidate extraction and multi-class rotated NMS expand tail latency ($\text{P95} = 80.85\,\text{ms}$, $\text{P99} = 95.91\,\text{ms}$), where PointPillars exhibits tighter tail dispersion ($75.88\,\text{ms}$ and $92.73\,\text{ms}$). Both architectures comfortably fulfill the strict $100\,\text{ms}$ automotive real-time constraint across $\ge 99\%$ of operational driving frames.

\textit{2. Component Latency Profiling and Host-to-Device Efficiency.}
Stage-level profiling identified the breakdown across pipeline stages, totaling $60.67\,\text{ms}$ across the four core modules plus a minor $0.29\,\text{ms}$ ($0.48\%$) evaluation harness residual (accounting for frame metadata indexing, tensor unbinding, and bounding-box formatting):
\begin{itemize}
    \item \textbf{Input Preprocessing ($20.94\,\text{ms}$)}: Point-cloud file loading and Strata-BEV14 rasterization with the Numba CPU backend took $20.94\,\text{ms}$ combined. The implementation uses occupied-cell accumulators and precomputed logarithmic lookup tables. Recorded preprocessing took $27.75\,\text{ms}$ for PointPillars and $27.20\,\text{ms}$ for SECOND. These measurements cover complete input preparation from disk to memory.
    \item \textbf{Host-to-Device Memory Transfer ($6.91\,\text{ms}$)}: Binary occupancy tensors with a 35-channel layout require $78.85\,\text{MB}$ per scan ($800 \times 704 \times 35 \times 4$ bytes). The Strata-BEV14 14-channel layout requires $31.54\,\text{MB}$, reducing tensor memory payload by exactly $2.50\times$ ($60.0\%$). Recorded H2D transfer time on Tesla T4 with pinned host memory buffers was $6.91\,\text{ms}$ for Strata-BEV14 ($31.54\,\text{MB}$), compared with $40.63\,\text{ms}$ for unpinned 35-channel transfers ($78.85\,\text{MB}$, Table~\ref{tab:encoding}). The resulting $5.88\times$ transfer acceleration reflects the combined effect of the $2.50\times$ payload reduction and page-locked pinned memory direct DMA transfers across the PCIe bus, which bypasses OS page paging.
    \item \textbf{Model Forward Propagation ($17.88\,\text{ms}$)}: Network inference within the sequential evaluation pipeline took $17.88\,\text{ms}$ ($55.94\,\text{FPS}$) on Tesla T4, outperforming MobilePIXOR's corresponding forward inference ($25.43\,\text{ms}$, $39.32\,\text{FPS}$). In the isolated benchmarking setup with pre-allocated tensors and \texttt{torch.cuda.Event} timing (Table~\ref{tab:components}), standalone model forward execution reached \textbf{115.32\,FPS} ($8.67\,\text{ms}$).
    \item \textbf{Decoding and Rotated NMS ($14.94\,\text{ms}$)}: Candidate peak selection and class-wise rotated NMS took $14.94\,\text{ms}$. Data remained on the GPU until final bounding box export.
\end{itemize}


\subsection{Distance-Stratified Detection Performance}
\label{sec:distance_stratified}



\paragraph{Range-Dependent Accuracy on Vulnerable Road Users.}
Because LiDAR beam divergence causes point density to drop quadratically over distance ($N \propto 1/r^2$), detector reliability degrades rapidly at intermediate and long ranges. Table~\ref{tab:distance_ap} benchmarks Moderate AP across three discrete distance intervals: Near ($0\text{--}30\,\text{m}$), Intermediate ($30\text{--}50\,\text{m}$), and Far ($50\text{--}70.4\,\text{m}$).

\begin{table}[!htbp]
\centering\footnotesize
\caption{Distance-stratified detection performance (Moderate AP, \%) on Vulnerable Road Users (Pedestrian and Cyclist) across Near ($0\text{--}30\,\text{m}$), Intermediate ($30\text{--}50\,\text{m}$), and Far ($50\text{--}70.4\,\text{m}$) range intervals on the KITTI validation benchmark using BEV AP $R_{40}$. Best results across competing compact models are in \textbf{bold}.}
\label{tab:distance_ap}
\setlength{\tabcolsep}{1.5pt}
\begin{tabular*}{\columnwidth}{@{\extracolsep{\fill}}llrrr@{}}
\toprule
\textbf{Category} & \textbf{Model} & \textbf{\shortstack{0--30\,m}} & \textbf{\shortstack{30--50\,m}} & \textbf{\shortstack{50--70.4\,m}} \\
\midrule
Pedestrian & MobilePIXOR~\cite{nguyen2024indoor} & 30.45 & 3.75 & \textbf{0.99} \\
& \textbf{Ours} & \textbf{65.23} & \textbf{11.49} & 0.80 \\
& \textit{Diff.} & \textit{+34.78} & \textit{+7.74} & \textit{-0.19} \\
\midrule
Cyclist & MobilePIXOR~\cite{nguyen2024indoor} & 25.55 & 12.61 & 0.56 \\
& \textbf{Ours} & \textbf{76.32} & \textbf{44.30} & \textbf{14.94} \\
& \textit{Diff.} & \textit{+50.77} & \textit{+31.69} & \textit{+14.38} \\
\bottomrule
\end{tabular*}
\end{table}




\textit{1. Near-Range Safety Envelope ($0\text{--}30\,\text{m}$).}
Compared with MobilePIXOR in the near-range $0\text{--}30\,\text{m}$ zone, MobilePIXORNeXt Pedestrian AP reached \textbf{65.23\%}, more than double MobilePIXOR's $30.45\%$ (a $+34.78\,\text{pp}$ improvement). Cyclist AP reached \textbf{76.32\%}, nearly triple MobilePIXOR's $25.55\%$ (a $+50.77\,\text{pp}$ improvement). These gains are consistent with Strata-BEV14's stratified height bins and continuous subcell offsets preserving fine boundaries when object returns are dense.

\textit{2. Intermediate and Long Ranges ($30\text{--}70.4\,\text{m}$).}
At medium ranges ($30\text{--}50\,\text{m}$), MobilePIXORNeXt achieved $11.49\%$ Pedestrian AP (vs. $3.75\%$) and $44.30\%$ Cyclist AP (vs. $12.61\%$, a $+31.69\,\text{pp}$ gain). Far-range ($50\text{--}70.4\,\text{m}$) Cyclist AP was $14.94\%$ (vs. $0.56\%$). These results are consistent with range-adaptive augmentation (\texttt{transform\_object}) and the wider context of $7 \times 7$ depthwise receptive fields helping under sparse returns. As qualitatively corroborated in our comparative visual analysis (Section~\ref{sec:qualitative_analysis}, Figure~\ref{fig:comparison_heatmap}), MobilePIXOR suffers severe activation attenuation on distant sparse returns, whereas MobilePIXORNeXt maintains focused focal peaks and confident bounding footprints.

\begin{table*}[t]
\centering\footnotesize
\caption{Source-linked inventory of historical configurations and the final detector. Labels A0--A10 are retained for navigation; they do not denote a sequence of experiments changing one factor at a time. AP is local BEV $R_{40}$ on KITTI validation. The schedule column records archived/completed epochs (planned total when incomplete), followed by selected checkpoint epoch.}
\label{tab:components}
\setlength{\tabcolsep}{2.5pt}
\begin{tabular*}{\textwidth}{@{\extracolsep{\fill}}ll>{\raggedright\arraybackslash}p{0.31\textwidth}rrrrr@{}}
\toprule
\textbf{Label} & \textbf{Run} & \textbf{Configuration} & \textbf{\shortstack{Epochs;\\selected}} & \textbf{\shortstack{mAP M\\(\%)}} & \textbf{\shortstack{mAP$_9$\\(\%)}} & \textbf{\shortstack{Model\\(ms)}} & \textbf{FPS} \\
\midrule
A0 & B0 & MobilePIXOR / Binary35 / legacy head & 38/100; 30 & 41.40 & 42.27 & 11.05 & 90.50 \\
A1 & B1 & Binary35 / single H16 / F24 SumFPN & 100; 13 & 47.60 & 49.08 & 5.88 & 169.97 \\
A2 & E8 & Strata-BEV8 / single H16 / F24 & 50; 20 & 47.37 & 48.95 & 6.06 & 165.02 \\
A3 & E14 & Strata-BEV14 / single H16 / F24 & 50; 30 & 50.14 & 52.00 & 5.92 & 169.01 \\
A4 & G16 & Grouped H16 / global augmentation / F24 & 50; 30 & 51.30 & 53.96 & 6.92 & 144.58 \\
A5 & O32 & Grouped H32 / simplified GT V1 / F32 SumFPN & 50; 50 & 60.22 & 61.11 & 6.70 & 149.15 \\
A6 & O32G16 & Grouped H16 / simplified GT V2 / F32 SG-FPN & 50; 50 & 66.25 & 68.45 & 7.45 & 134.25 \\
A7 & H32 & Grouped H32 / hybrid GT (6/8/10) / SumFPN & 50; 50 & 68.12 & 70.05 & 6.72 & 148.82 \\
A8 & S32 & Grouped H32 / hybrid GT (15/10/10) / SG-FPN & 50; 50 & 68.46 & 70.28 & 6.94 & 144.16 \\
A8F & P1 & Focal context / baseline loss / no IQA & 50; 50 & 67.38 & 69.46 & 7.77 & 128.66 \\
A9 & P4 & Focal context / OGA + IQA / independent towers & 50; 50 & 70.58 & 72.66 & 8.19 & 122.16 \\
A10 & Final S42 & No focal / Block-SimAM / shared grouped towers & 160; 130 & 71.50 & 73.45 & 8.67 & 115.32 \\
\bottomrule
\end{tabular*}
\par\vspace{2pt}
{\footnotesize \textit{Note.} B1 uses loss-based selection; the other listed AP checkpoints follow their archived selection records. A5 changes both head width and sampler relative to A4. A6 is a simplified-sampler H16 run, rather than the claimed PC-GTA H32/F24 addition. A7--A8 also change insertion quotas. A8F identifies the focal-context P1 result previously conflated with A8. A9 and A10 have different context modules, head architectures, and schedules. Historical model timings are archived values; the final model timing is from L4. These rows do not isolate backbone, head, PC-GTA, neck, or Block-SimAM gains under a common budget.}
\end{table*}


\subsection{Ablation Studies}
\label{sec:ablations}

\paragraph{Systematic Architectural Progression.}
To quantify the individual and cumulative contributions of the proposed components, Table~\ref{tab:components} details the step-by-step ablation progression from the baseline to the final configuration.

\paragraph{1. Backbone Architecture (A0 $\to$ A1).}
Replacing MobilePIXOR's residual blocks with a ConvNeXt-style backbone using $7 \times 7$ depthwise convolutions, inverted expansion $e=2.5$, SiLU, and LayerScale raised Moderate mAP from $41.40\%$ to $47.60\%$ (\textbf{+6.20\,pp Moderate mAP}), while standalone throughput nearly doubled from $90.57$ to \textbf{169.97\,FPS}. Large depthwise kernels broaden BEV receptive fields with limited computational cost.

\paragraph{2. BEV Input Representation (A1 $\to$ A2 $\to$ A3).}
\begin{table*}[t]
\centering\small
\caption{Archived BEV representation results on KITTI validation using the local evaluator with $R_{40}$ interpolation. The early configurations use a MobilePIXORNeXt backbone, F24 neck, and single H16 head. Training budgets and checkpoint-selection criteria differ, so these are descriptive representation comparisons. Payload is the FP32 tensor size for an $800\times704$ grid.}
\label{tab:encoding}
\setlength{\tabcolsep}{2.5pt}
\begin{tabular*}{\textwidth}{@{\extracolsep{\fill}}lcrrrrrrr@{}}
\toprule
\textbf{Representation} & \textbf{$C_{\mathrm{in}}$} & \textbf{\shortstack{H2D Payload\\(MB) $\downarrow$}} & \textbf{\shortstack{Car AP$_{\mathrm{M}}$\\(\%) $\uparrow$}} & \textbf{\shortstack{Pedestrian\\AP$_{\mathrm{M}}$ (\%) $\uparrow$}} & \textbf{\shortstack{Cyclist\\AP$_{\mathrm{M}}$ (\%) $\uparrow$}} & \textbf{\shortstack{mAP$_{\mathrm{M}}$\\(\%) $\uparrow$}} & \textbf{\shortstack{mAP$_{9}$\\(\%) $\uparrow$}} & \textbf{\shortstack{Archived model\\FPS $\uparrow$}} \\
\midrule
TriBand-BEV~\cite{khoshkdahan2026triband} & 3 & 6.76 & 74.10 & 38.05 & 21.12 & 44.42 & 46.55 & 188.43 \\
Strata-BEV8  & 8 & 18.02 & 78.35 & 38.48 & 25.29 & 47.37 & 48.95 & 165.02 \\
Binary35~\cite{yang2018pixor} & 35 & 78.85 & 75.24 & 44.36 & 23.20 & 47.60 & 49.08 & 169.97 \\
\textbf{Strata-BEV14 (Ours)} & 14 & 31.54 & 79.97 & 45.36 & 25.10 & 50.14 & 52.00 & 169.01 \\
\bottomrule
\end{tabular*}
\par\vspace{2pt}
{\footnotesize \textit{Note.} Binary35 is the historical B1 run (100-epoch budget, loss-selected epoch 13); Strata-BEV8 and Strata-BEV14 are E8 and E14 (50-epoch budgets, AP-selected epochs 20 and 30). The reported forward rates are archived measurements, rather than a new common-T4 benchmark. The 14-channel tensor is 60.0\% smaller than Binary35. The AP gap does not isolate the effects of individual height statistics or centroid offsets. TriBand-BEV denotes the local three-channel adaptation used in this experiment, rather than a reproduction of the complete cited detector.}
\end{table*}

To assess how input rasterization affects detection accuracy and latency, Table~\ref{tab:encoding} compares four distinct BEV representations under identical MobilePIXORNeXt backbone and single-head architecture. TriBand-BEV~\cite{khoshkdahan2026triband} partitions points into three coarse height bands; while it delivers the highest inference throughput ($188.4\,\text{FPS}$) and lowest Host-to-Device transfer payload ($6.76\,\text{MB}$), its Moderate mAP drops to $44.42\%$ ($-5.72\,\text{pp}$ below Strata-BEV14). The lack of continuous elevation moments and horizontal subcell centroid offsets degrades small object localization (Pedestrian AP $38.05\%$, Cyclist AP $21.12\%$). Conventional 35-channel binary occupancy slicing (A1) achieves $47.60\%$ mAP but requires transferring $78.85\,\text{MB}$ per scan, imposing significant Host-to-Device transfer latency ($40.63\,\text{ms}$) while still discarding sub-voxel point coordinates. The 8-channel Strata-BEV8 representation (A2) yields comparable accuracy ($47.37\%$) with fewer channels. In contrast, our 14-channel Strata-BEV14 representation (A3) achieves the highest Moderate mAP of $50.14\%$ ($+5.72\,\text{pp}$ over TriBand-BEV and $+2.54\,\text{pp}$ over Binary35) at $169.0\,\text{FPS}$. Continuous vertical spatial moments ($z_{\max}, z_{\mathrm{mean}}, z_{\mathrm{span}}, z_{\mathrm{std}}$) and subcell centroid displacement offsets $(\Delta\bar{x}, \Delta\bar{y})$ retain critical 3D geometry otherwise lost through horizontal discretization without incurring PointNet encoder overhead.

\paragraph{3. Decoupled Grouped Prediction Towers (A3 $\to$ A4 $\to$ A5).}
Separating Car and Ped/Cyc prediction heads (A4) raised Moderate mAP to $51.30\%$ ($+1.16\,\text{pp}$ over A3). Increasing grouped head width from 16 to 32 channels (A5) raised Moderate mAP to $60.22\%$ (\textbf{+8.92\,pp Moderate mAP}) and $\text{mAP}_9$ to $61.11\%$ relative to configuration A4 (a combined $+10.08\,\text{pp}$ gain over A3). In configuration A4, compressing shared intermediate features into only 16 channels induces severe subspace rank collapse when the convolutional tower attempts to project four distinct regression targets (offset, log-size, doubled-angle yaw, and quality) across heterogeneous object scales. Expanding head width to 32 channels (A5) adds only $75\,\text{K}$ parameters but doubles the latent representation rank. This additional subspace capacity provides sufficient degrees of freedom to decouple inter-task gradients between orthogonal geometric properties, driving an $+8.92\,\text{pp}$ jump ($+10.08\,\text{pp}$ over baseline A3).

\paragraph{4. Physics-Consistent Ground-Truth Augmentation (A5 $\to$ A6).}
Physics-consistent augmentation (\texttt{PC-GTA}) improved Moderate mAP by \textbf{+6.03\,pp} over A5, from $60.22\%$ to $66.25\%$. It also raised $\text{mAP}_9$ from $61.11\%$ to $68.45\%$ (\textbf{+7.34\,pp $\text{mAP}_9$}). The evaluated sampler combined terrain alignment, static obstacle rejection, and range-dependent point subsampling. Optional LOS, shadow-masking, and radiometric mechanisms were disabled in the final run.

\paragraph{5. Multiscale Neck and Scale Gating (A6 $\to$ A7 $\to$ A8).}
Increasing neck width from 24 to 32 channels (A7) raised Moderate mAP to $68.12\%$ (\textbf{+1.87\,pp Moderate mAP}). The Top-Down Scale-Gated FPN (A8) replaced transposed convolutions with bilinear upsampling and learnable depthwise gating to avoid deconvolution checkerboard artifacts in continuous regression maps. Under standard unweighted Smooth-L1 coordinate regression, this depthwise gating causes a slight drop from $68.12\%$ to $67.38\%$ ($-0.74\,\text{pp}$), as multiplicative gating without explicit geometric loss coupling tends to attenuate weak gradient signals from distant and sparse point returns.

\paragraph{6. Geometric Supervision and Localization Quality (A8 $\to$ A9).}
Oriented Geometric Alignment (OGA) and IoU-Aware Quality Alignment (IQA) raised Moderate mAP from $67.38\%$ to $70.58\%$ (\textbf{+3.20\,pp Moderate mAP}). This $+3.20\,\text{pp}$ surge reveals a decisive co-design synergy: while SG-FPN alone slightly attenuates weak returns under Smooth-L1, it provides the precise multi-scale spatial separation that OGA and IQA require. Multi-Axis Projection GIoU and scale-normalized $\pi$-symmetric corner distances provide direct, rotation-aware gradient feedback across the scale-gated features, while IQA calibrates classification confidence against continuous projection-based geometric quality ($\hat{q}^*$) to eliminate the suppression of accurate candidates by poorly aligned boxes during rotated NMS.

\paragraph{7. Block-SimAM Energy Attention (A9 $\to$ A10).}
Embedding parameter-free Block-SimAM inside all nine residual blocks across stages C2, C3, and C4 (depths $3, 4, 2$) produced the final configuration A10, reaching \textbf{71.50\% Moderate mAP} and \textbf{73.45\% mAP$_9$} at the converged 100-epoch deployment checkpoint (\textbf{+0.92\,pp Moderate mAP} over A9, while reaching $71.48\%$ and $73.57\%$ at 50 epochs).
Crucially, because both A9 and A10 maintain the identical backbone depth of nine blocks ($3+4+2$) and identical head structures, this $+0.92\,\text{pp}$ accuracy gain is strictly attributable to 3D spatial-energy attention without any confounding depth alterations or learnable parameter inflation (adding exactly 0 parameters). The minimal increase of $0.48\,\text{ms}$ in standalone forward latency ($8.19 \to 8.67\,\text{ms}$) confirms the extreme efficiency of the analytical energy minimization.
Block-SimAM gained 1.66--2.10\,pp in Moderate mAP over stage-level variants at 115.32 standalone FPS (versus 144.05\,FPS for Stage-3 SimAM).


\paragraph{Geometric Supervision and Quality Head Factorial Study.}
Table~\ref{tab:loss} presents a $2 \times 2$ factorial evaluation isolating the effects of OGA geometric supervision and the explicit IQA prediction branch.

\begin{table*}[t]
\centering\small
\caption{Geometric supervision (Baseline Smooth-L1 vs. OGA), localization quality prediction (with and without an IQA Head), and head-free Q-OGA on KITTI validation at Moderate difficulty ($R_{40}$). IQA models rank NMS candidates using $S_{\mathrm{det}} = \sqrt{P_{\mathrm{cls}} \hat{q}}$. Standalone forward latency is measured in FP32 on an NVIDIA Tesla T4 GPU. Best values are in \textbf{bold}, and second-best are \underline{underlined}.}
\label{tab:loss}
\setlength{\tabcolsep}{3pt}
\begin{tabular*}{\textwidth}{@{\extracolsep{\fill}}>{\raggedright\arraybackslash}p{0.26\textwidth}ccrrrrr@{}}
\toprule
\textbf{Training Objective} & \textbf{\shortstack{IQA\\Head}} & \textbf{\shortstack{Params\\(M) $\downarrow$}} & \textbf{\shortstack{Forward\\(ms) $\downarrow$}} & \textbf{\shortstack{Car AP$_{\mathrm{M}}$\\(\%) $\uparrow$}} & \textbf{\shortstack{Pedestrian\\AP$_{\mathrm{M}}$ (\%) $\uparrow$}} & \textbf{\shortstack{Cyclist\\AP$_{\mathrm{M}}$ (\%) $\uparrow$}} & \textbf{\shortstack{mAP$_{\mathrm{M}}$\\(\%) $\uparrow$}} \\
\midrule
Baseline Smooth-L1 & No  & \textbf{0.702} & \textbf{6.94} & 83.92 & 55.48 & 62.74 & 67.38 \\
OGA (MGIoU + $\pi$-Corner) & No  & \textbf{0.702} & \textbf{6.94} & 84.09 & 57.93 & 64.88 & 68.97 \\
Baseline Smooth-L1 & Yes & 0.721 & 8.19 & 84.13 & \underline{59.45} & \textbf{67.34} & 70.31 \\
\textbf{OGA (MGIoU + $\pi$-Corner)} & \textbf{Yes} & 0.721 & 8.19 & \textbf{84.87} & \textbf{59.63} & \underline{67.24} & \textbf{70.58} \\
\midrule
Q-OGA (Quality-Coupled Focal Loss) & No & \textbf{0.702} & \textbf{6.94} & \underline{84.35} & 58.12 & 65.13 & 69.20 \\
\bottomrule
\end{tabular*}
\par\vspace{2pt}
{\footnotesize \textit{Note.} All variants share identical MobilePIXORNeXt baseline backbone with legacy stage depths $(2, 4, 2)$, Strata-BEV14 input, Grouped H32 heads, F32 SG-FPN neck, and PC-GTA augmentation ($0.721\,\text{M}$ parameters with IQA, $0.702\,\text{M}$ without IQA). The final consolidated configuration A10 ($[3, 4, 2]$) contains $753{,}569$ parameters ($0.754\,\text{M}$), confirming that Block-SimAM adds zero parameters. When IQA is enabled, candidate boxes are ranked during NMS via composite score $S_{\mathrm{det}} = \sqrt{P_{\mathrm{cls}} \cdot \hat{q}}$. Q-OGA integrates quality directly into classification supervision without dedicating a separate head, saving $18{,}626$ parameters ($18.63\,\text{K}$) and $1.25\,\text{ms}$ forward latency compared with explicit IQA models.}
\end{table*}


Table~\ref{tab:loss} reports the OGA and IQA ablation. Notably, in the absence of an auxiliary IQA branch, OGA alone raises Moderate mAP from $67.38\%$ to $68.97\%$ (\textbf{+1.59\,pp Moderate mAP} over baseline Smooth-$L_1$), with decisive gains across all three classes: Car ($+0.17\,\text{pp}$), Pedestrian ($+2.45\,\text{pp}$), and Cyclist ($+2.14\,\text{pp}$). This confirms that rotation-aware geometric alignment independently guides coordinate convergence even without confidence re-ranking. When IQA is enabled, OGA further elevates Car AP to its benchmark peak of $\mathbf{84.87\%}$ (versus $84.13\%$ for Baseline + IQA), yielding a joint optimum of \textbf{70.58\% Moderate mAP}. These results demonstrate clear methodological synergy: MGIoU and scale-normalized $\pi$-corner distances enforce strict boundary and orientation gradients during backpropagation, while IQA adjusts detection rank via $S_{\mathrm{det}} = \sqrt{P_{\mathrm{cls}}\hat{q}}$ to prioritize geometrically faithful candidates during rotated NMS.

Table~\ref{tab:loss} compares explicit quality prediction with head-free quality coupling. IQA models contained $0.754\,\text{M}$ parameters ($753{,}569$ parameters). Q-OGA (Eq.~(\ref{eq:cqfl_target})--(\ref{eq:cqfl_loss})) omitted the IQA heads, using $0.735\,\text{M}$ parameters ($734{,}943$ parameters, saving $18{,}626$ parameters). Without extra prediction branches or post-processing score multiplication, Q-OGA achieved $69.20\%$ Moderate mAP, exceeding OGA alone by $+0.23\,\text{pp}$ and baseline Smooth-L1 by $+1.82\,\text{pp}$. Q-OGA therefore offers a smaller, lightweight configuration for constrained deployment, while explicit IQA achieved the highest accuracy ($70.58\%$) with a dedicated quality branch.

\paragraph{3D Spatial-Energy Attention Integration and Placement.}
Table~\ref{tab:attention} investigates the placement of parameter-free SimAM energy attention.

\begin{table*}[t]
\centering\normalsize
\caption{Historical attention-related configurations and final-detector snapshot on KITTI validation (local BEV AP $R_{40}$). Neck width, focal context, prediction towers, schedule, and checkpoint identity are not held fixed across all rows; this table is not a controlled placement ablation.}
\label{tab:attention}
\setlength{\tabcolsep}{4pt}
\begin{tabular}{@{}llrrrr@{}}
\toprule
\textbf{\shortstack[l]{SimAM Attention\\Placement}} & \textbf{\shortstack[l]{Neck Fusion\\Width}} & \textbf{\shortstack{mAP$_{\mathrm{M}}$\\(\%) $\uparrow$}} & \textbf{\shortstack{mAP$_{9}$\\(\%) $\uparrow$}} & \textbf{\shortstack{Archived model\\Forward (ms) $\downarrow$}} & \textbf{\shortstack{Archived model\\FPS $\uparrow$}} \\
\midrule
None (P4, focal context) & $F_{32}$ & 70.58 & 72.66 & 8.19 & 122.16 \\
Stage-2 ($C_2$) & $F_{32}$ & 69.44 & 71.49 & 7.30 & 136.91 \\
Stage-3 ($C_3$, R25) & $F_{24}$ & 69.80 & 71.98 & 6.94 & 144.05 \\
Stage-4 ($C_4$) & $F_{32}$ & 69.36 & 71.34 & 7.44 & 134.43 \\
\textbf{Block-level ($C_2, C_3, C_4$)} & \textbf{$F_{32}$} & 71.50 & 73.45 & 8.67 & 115.32 \\
\bottomrule
\end{tabular}
\par\vspace{2pt}
{\footnotesize \textit{Note.} The no-attention row is historical P4 (852,929 parameters, focal context); the final Block-SimAM row is the seed-42 shared-head detector (753,569 parameters, no focal context). The archived R25 stage-attention row uses F24, unlike the F32 final detector. The final 8.67\,ms model time comes from the L4 report; a common GPU/timing boundary has not been established for every row. Accordingly, differences in AP or latency cannot be assigned solely to attention placement. The separate T4 pipeline timing is reported in Table~\ref{tab:efficiency}.}
\end{table*}

Stage-level SimAM at downsampling boundaries ($C_2, C_3, C_4$) reduced mAP to $69.44\%$, $69.80\%$, and $69.36\%$, below the no-attention score ($70.58\%$). Attention at these transitions may blur fine spatial gradients after downsampling. Block-SimAM inside every residual block instead achieved \textbf{71.50\% Moderate mAP} (\textbf{+0.92\,pp gain} over no attention) and \textbf{73.45\% mAP$_9$} at the converged 100-epoch deployment checkpoint ($71.48\%$ and $73.57\%$ at 50 epochs). It modulates the residual update before LayerScale addition ($\mathbf{Z}' = \mathbf{Z} + \boldsymbol{\gamma}_{\mathrm{blk}} \odot \mathcal{S}(\mathbf{R})$) without changing inter-stage feature resolution, allowing modulation of sparse returns within each block.

\paragraph{Multi-Task Optimization Dynamics and Convergence.}
Figure~\ref{fig:loss_curves} plots training and validation losses over 50 epochs for configurations A1 through A10. It includes the total objective and six tasks: Classification, Geometry, IoU Quality, Subcell Offset, Size, and Yaw. We used these trajectories to assess training stability and convergence.
Across 50 epochs, total and task losses exhibited steady downward convergence trends. Total loss fell from $10^4$ to approximately $1.0$ within the first 15 epochs, progressing stably through cosine annealing. While transient validation classification spikes were observed between epochs 18 and 25 (briefly reaching $\sim 10^2$ before dampening), the training and validation regression curves remained well-bounded without unbounded divergence. This trend reflects the stabilizing effect of physics-consistent data augmentation (PC-GTA, A6) and loss balancing.
Temperature-Softmax Bounded Uncertainty Weighting (T-SBUW, Eq.~(\ref{eq:tsbuw_weights})--(\ref{eq:balanced_loss})) dynamically bounded task loss weights ($\omega_k \in (0, J)$), mitigating magnitude disparities and objective competition between cross-entropy classification and multi-task coordinate regression. Configuration A10 achieved low, stable late-training validation losses across regression tasks, aligning with the empirical evaluation in Table~\ref{tab:components}.

\paragraph{Training Schedule Fairness and Convergence Validation.}
To evaluate training duration and convergence dynamics within MobilePIXORNeXt, we benchmark the rapid ablation checkpoint from the 50-epoch schedule (Epoch~47) against the fully converged 100-epoch deployment checkpoint (Epoch~100). Under the local loader-aligned BEV evaluator with $R_{40}$ interpolation, the 50-epoch checkpoint attained $71.48\%$ Moderate mAP ($82.71\%$ Car, $62.30\%$ Pedestrian, $69.43\%$ Cyclist, and $73.57\%$ $\text{mAP}_9$), closely matching the $71.50\%$ Moderate mAP and $73.45\%$ $\text{mAP}_9$ achieved by the 100-epoch deployment model ($86.90\%$ Car, $60.78\%$ Pedestrian, $66.81\%$ Cyclist). These consistent aggregate figures confirm that the rapid 50-epoch schedule used for the architectural ablations (Table~\ref{tab:components}) faithfully captures converged model performance without distortion. The essential benefit of training to 100 epochs is the progressive calibration of orientation and sharpening of foreground classification heatmaps (Figure~\ref{fig:comparison_heatmap}), which systematically suppresses low-confidence background false alarms. By drastically reducing the candidate proposals routed to post-processing and rotated NMS, this convergence effect slashes end-to-end pipeline latency by $25.0\%$ ($81.23\,\text{ms} \to 60.96\,\text{ms}$, elevating throughput from $12.31$ to $16.41\,\text{FPS}$).


\begin{figure*}[p]
\centering
\paperfigure[width=\textwidth,height=0.88\textheight,keepaspectratio]{training_curves_all_losses.pdf}{0.88\textheight}{Training and validation loss trajectories across 50 epochs.}
\captionof{figure}{Empirical multi-task loss trajectories over 50 training epochs on the KITTI benchmark for architectural configurations A1--A10 (Table~\ref{tab:components}). Subpanels depict: \textbf{(a)} Total training loss, \textbf{(b)} Total validation loss, \textbf{(c)} Focal classification loss, \textbf{(d)} Geometric OGA loss ($\ell_{\mathrm{proj}} + \ell_{\mathrm{corner}}$), \textbf{(e)} IQA quality alignment loss, \textbf{(f)} Subcell horizontal offset loss, \textbf{(g)} Footprint bounding dimension loss, and \textbf{(h)} Doubled-angle yaw orientation loss. Trajectories illustrate stable monotonic descent from initial magnitudes ($10^4$) to asymptotic convergence ($\sim 1.0$), with transient validation classification fluctuations at epochs 18--25 rapidly stabilized by T-SBUW adaptive weighting.}
\label{fig:loss_curves}
\end{figure*}



\subsection{Qualitative Visualizations and Analysis}
\label{sec:qualitative_analysis}

\begin{figure*}[!t]
\centering
\paperfigure[width=\textwidth]{comparison_mobilepixor_heatmap.pdf}{65mm}{Comparative BEV detection and classification heatmap visualizations between MobilePIXOR and MobilePIXORNeXt.}
\captionof{figure}{Qualitative comparison of Bird's-Eye View (BEV) detections and predicted classification heatmaps between the baseline MobilePIXOR and MobilePIXORNeXt (Ours) across two representative driving scenes. \textbf{Row A}: MobilePIXOR exhibits diffused, low-amplitude heatmap responses with noticeable background clutter, leading to missed detections (false negatives) on distant and sparse point clusters. \textbf{Row B}: MobilePIXORNeXt produces sharp, concentrated focal peaks strictly centered on vehicle targets with suppressed background activations, successfully recovering distant objects and yielding tightly aligned oriented bounding footprints.}
\label{fig:comparison_heatmap}
\end{figure*}

\paragraph{Comparative Heatmap and Detection Analysis against Baseline.}
Figure~\ref{fig:comparison_heatmap} provides a side-by-side qualitative comparison between the baseline MobilePIXOR and MobilePIXORNeXt across two challenging driving scenarios, illustrating both oriented BEV bounding footprints and the corresponding foreground classification heatmaps.

Two critical observations emerge from this comparison:
\begin{itemize}
    \item \textbf{Focal Heatmap Concentration and Clutter Suppression}: In MobilePIXOR (Row~A), the classification heatmaps exhibit diffuse, low-intensity activation clusters (frequently below $0.35$ normalized confidence) accompanied by spurious background activations in uninformative regions. In contrast, MobilePIXORNeXt (Row~B) generates sharp, high-magnitude focal peaks ($>0.80$) centered precisely at vehicle centroids, while effectively suppressing background clutter to near-zero activation. This sharpening directly validates the efficacy of our Block-SimAM spatial-energy attention and Top-Down Scale-Gated FPN, which collaboratively amplify sparse point features while filtering scale-inconsistent lateral noise. As analyzed in Section~\ref{sec:ablations}, this background suppression reduces low-confidence candidate proposals, directly speeding up Rotated NMS.
    \item \textbf{Distant and Sparse Object Recall}: At intermediate and long ranges ($>30\,\text{m}$), where LiDAR point returns become severely attenuated ($N \propto 1/r^2$), MobilePIXOR fails to exceed the peak detection threshold, producing false negatives. MobilePIXORNeXt reliably recovers these distant vehicles with confident heatmap activations and well-oriented bounding boxes. This visual evidence corroborates the substantial $+31.69\,\text{pp}$ Cyclist AP and $+7.74\,\text{pp}$ Pedestrian AP gains reported in Table~\ref{tab:distance_ap} for the $30\text{--}50\,\text{m}$ interval, demonstrating the benefits of Strata-BEV14's continuous elevation moments and the enlarged $7 \times 7$ depthwise receptive field.
\end{itemize}

\paragraph{Visual Inspection on Representative Driving Scenarios.}
Figure~\ref{fig:qualitative_detections} presents MobilePIXORNeXt predictions on three KITTI validation frames (top to bottom: \texttt{007038}, \texttt{005571}, and \texttt{006761}). Class labels and confidence scores are retained in the camera overlays to support inspection of the predicted footprints; the BEV panels retain class annotations.

\begin{figure*}[!t]
\centering
\paperfigure[width=\textwidth,height=0.38\textheight,keepaspectratio]{qualitative_detections.pdf}{65mm}{Qualitative BEV detection visualizations on KITTI validation scenes.}
\captionof{figure}{Qualitative BEV predictions from MobilePIXORNeXt (converged MobilePIXORNeXt deployment model) on KITTI validation frames \texttt{007038}, \texttt{005571}, and \texttt{006761} (top to bottom). \textbf{Left}: Front-camera images with calibrated projections of predicted BEV footprints, class labels, and confidence scores. \textbf{Right}: Corresponding LiDAR BEV views with predicted oriented footprints and class annotations. Prediction colors are cyan for Car, magenta for Pedestrian, and yellow for Cyclist. The panels visualize predicted footprints rather than full 3D boxes.}
\label{fig:qualitative_detections}
\end{figure*}

The paired camera and BEV views in Figure~\ref{fig:qualitative_detections} support inspection of the following prediction patterns:
\begin{itemize}
    \item \textbf{Vehicle Predictions}: Frame \texttt{007038} contains several predicted vehicle footprints at different forward distances, visible in both the camera projection and the BEV panel.
    \item \textbf{Mixed-Class Predictions}: Frame \texttt{005571} contains Car, Pedestrian, and Cyclist predictions. The retained labels and confidence scores identify the classes assigned to nearby footprints.
    \item \textbf{Nearby Road Users}: Frame \texttt{006761} shows predicted footprints around vehicles, pedestrians, and cyclists in the same scene. The two views reveal the spatial arrangement of these predictions.
\end{itemize}

\subsection{Discussion and Practical Limitations}
\label{sec:discussion_limitations}

\paragraph{Discussion and Insights.}
The results showed that accurate 3D LiDAR detection was possible without large point-cloud encoders or 3D sparse convolutions. The deterministic 2D BEV representation retained multi-band densities, elevation dispersion, and continuous subcell centroid displacements. A ConvNeXt-style 2D network extracted geometric features with fewer than $0.76\,\text{M}$ parameters. Physics-consistent augmentation (\texttt{PC-GTA}) and decoupled grouped heads provided the most substantial individual accuracy improvements in the architectural progression, demonstrating the value of domain-specific LiDAR data modeling and multi-class gradient isolation.

\paragraph{Limitations and Future Scope.}
MobilePIXORNeXt is formulated specifically for high-throughput, oriented 2D Bird's-Eye View footprint detection on the $SE(2)$ manifold, directly fulfilling the horizontal clearance and Time-to-Collision (TTC) constraints of vehicular collision avoidance. While an auxiliary vertical regression branch can be optionally incorporated to recover complete 3D cuboid parameters $(z, h)$, our core benchmark focuses on $SE(2)$ BEV localization. Several operational boundaries remain. First, our empirical validation relies on 64-beam mechanical spinning LiDAR on predominantly planar urban roadways. Although Strata-BEV14's continuous elevation moments ($z_{\mathrm{span}}, z_{\mathrm{std}}$) and log-density normalization maintain numerical stability under point sparsity, performance on ultra-sparse 16-beam or non-repetitive solid-state LiDAR sensors warrants systematic domain investigation. Second, on severe terrain slopes or multi-level grade transitions, planar BEV projections introduce perspective distortion along vertical structural boundaries. Future work will address these constraints by developing Range-Conditioned Scale-Gated Feature Pyramids (RC-SGFPN) with continuous spatial frequency modulation, and deploying hardware-tailored TensorRT INT8 quantization for low-power edge microcontrollers.
